\documentclass[11pt,twoside,openright]{book}

\usepackage[T1]{fontenc}
\usepackage[utf8]{inputenc}
\usepackage{amsmath,amssymb,amsthm,mathtools,bm}
\usepackage{mathpazo}
\usepackage{microtype}
\usepackage[inner=1.35in,outer=1.1in,top=1.25in,bottom=1.35in]{geometry}
\usepackage{booktabs}
\usepackage{fancyhdr}
\usepackage[hidelinks]{hyperref}
\usepackage{cleveref}

\setlength{\parindent}{1.4em}
\setlength{\parskip}{0pt}
\linespread{1.07}
\emergencystretch=3em

\pagestyle{fancy}
\fancyhf{}
\fancyhead[LE]{\small\thepage\quad\textsc{A Cinema of Individual Editions}}
\fancyhead[RO]{\small\nouppercase{\leftmark}\quad\thepage}
\renewcommand{\headrulewidth}{0.3pt}
\fancypagestyle{plain}{\fancyhf{}\fancyfoot[C]{\small\thepage}\renewcommand{\headrulewidth}{0pt}}

\theoremstyle{plain}
\newtheorem{theorem}{Theorem}[chapter]
\newtheorem{proposition}[theorem]{Proposition}
\newtheorem{lemma}[theorem]{Lemma}
\newtheorem{corollary}[theorem]{Corollary}
\theoremstyle{definition}
\newtheorem{definition}[theorem]{Definition}
\newtheorem{example}[theorem]{Example}
\theoremstyle{remark}
\newtheorem{remark}[theorem]{Remark}

\newcommand{\Capac}{\mathrm{Cap}}
\newcommand{\Rob}{\mathrm{Rob}}
\newcommand{\Inv}{\mathrm{Inv}}
\DeclareMathOperator{\loss}{loss}

\title{\Huge\bfseries A Cinema of Individual Editions\\[1.2ex]
\large Chapter Fourteen: Sequence and Reinterpretation\\[0.4ex]
\large Working Draft}
\author{Flyxion\\[0.3ex]\small Independent Researcher}
\date{November 2026}

\begin{document}

\frontmatter
\maketitle
\tableofcontents
\mainmatter
\setcounter{chapter}{13}

\chapter{Sequence and Reinterpretation}
\label{ch:sequence}

The budget of Chapter~13 treats films as items. A viewer has so many viewings, a catalogue supplies so many productions, and the arithmetic is the same in whatever order the films are seen. Films are not items. A later film can change what an earlier one means, and what a film means at the first viewing is a fact about the viewer's life at that moment and cannot be had again. The order of release is therefore a variable, and a book about individual editions has to say what it varies.

This is the most open chapter of the book and it is written as a program. It states a model in which order can be seen to matter and in which it can be seen not to, computes what the model says on the engine scene of Chapter~3, and marks where the model stops and an experiment would have to begin. It does not claim that releasing a film in the tenth year of a viewer's life does something that releasing it in the first would miss. It states what would have to be true for that to hold.

\section{Films as Commitments}\label{sec:commit14}

The engine scene of Chapter~3 has $22$ histories and four readings: diligence, negligence, disclosure, and concealment. An edition, by the account of that chapter, commits to a set of histories, those that agree with everything it shows. Let $C_f\subseteq H$ be the set committed to by a film $f$, which in this chapter may be an edition of any production set in the same story world. A viewer's \emph{knowledge} $K\subseteq H$ is the set of histories not yet excluded by what the viewer has seen. The \emph{reading set} of a knowledge state is the set of readings of the histories in it, and its size is the \emph{openness}. A viewer who has seen nothing has all $22$ histories and all four readings.

A viewing updates knowledge. Two updates are considered. The \emph{intersection} update takes $K\cap C_f$ and is defined when that set is not empty. The \emph{revision} update takes $K\cap C_f$ when it is not empty and otherwise takes $C_f$, so that a film that contradicts everything the viewer believed replaces it and the latest film has priority. Neither is a theory of how people understand stories. They are the two simplest ways of letting a later film act on an earlier one, and they behave differently.

\begin{proposition}[Order and final state]\label{prop:order14}
For a family of films whose committed sets have a nonempty common intersection, every order of viewing leaves the viewer with the same knowledge, under the intersection update and under the revision update alike, namely $H\cap\bigcap_fC_f$. For two films whose committed sets are disjoint, the revision update leaves the viewer with the set of whichever film came last, and so the final state depends on the order.
\end{proposition}

\begin{proof}
If the common intersection is nonempty, every partial intersection is nonempty, so neither update ever takes its second branch, and intersection is commutative. If the two sets are disjoint and nonempty, the second film's set replaces the first in either order, the two final states are the two different sets, and they are different because disjoint nonempty sets are never equal.
\end{proof}

The proposition separates two things that a discussion of reinterpretation tends to run together. The final knowledge of a viewer depends on order only when films contradict one another. What can depend on order in the other cases is the path. The \emph{path} here is the sequence of knowledge states through which the viewer passed, and a summary of it, the openness of each film at its first viewing, is what the tables below report. Repeated or redundant films can give identical paths in different orders, so the claim is that paths can differ, and not that they differ in every order. The path is where the experience of a film lies.

\section{Retrospective Reinterpretation}\label{sec:retro14}

Take the $32$ faithful editions of the engine scene, those consistent with the script history, with the distribution of openness in \Cref{tab:open14}: eight leave one reading, ten leave two, eight leave three, and six leave all four, with a mean of $2.375$.

\begin{table}[ht]
\centering
\begin{tabular}{@{}lcccc@{}}
\toprule
Readings left open & 1 & 2 & 3 & 4\\
\midrule
Faithful editions & 8 & 10 & 8 & 6\\
\bottomrule
\end{tabular}
\caption{Openness of the $32$ faithful editions of the engine scene.}\label{tab:open14}
\end{table}

Now release a later film in the same world, and let the viewer apply it to what an earlier edition left open. \Cref{tab:retro14} gives four later films by what they commit to. The first shows that a character noticed the bolt, and commits to $14$ of the $22$ histories. The second shows that a character noticed and concealed, committing to seven. The third shows the third bolt loose, committing to twelve. The fourth is unrelated and commits to all. A fifth, in the last row, contradicts the script by showing the third bolt tight, committing to ten.

\begin{table}[ht]
\centering
\begin{tabular}{@{}lp{2.6cm}p{3cm}p{2.4cm}@{}}
\toprule
Later film & Reading set shrinks & Mean openness, before $\to$ after & Inconsistent with the edition\\
\midrule
Noticed & 16 of 32 & $2.375\to1.500$ & 0\\
Noticed and concealed & 24 of 32 & $2.375\to1.000$ & 0\\
Third bolt loose & 12 of 32 & $2.375\to2.000$ & 0\\
Unrelated & 0 of 32 & $2.375\to2.375$ & 0\\
Third bolt tight & 0 of 24 & $2.500\to2.500$ & 8\\
\bottomrule
\end{tabular}
\caption{Effect of one later film on the reading sets of the faithful editions, by the intersection update. In the last row the eight editions that commit to the third bolt loose contradict the later film, and the means are over the other twenty-four.}\label{tab:retro14}
\end{table}

A later film that commits to the concealment settles twenty-four of the thirty-two earlier editions. A later film that does not bear on the question leaves every reading set where it was. No later film in the table widens any reading set, and this is general.

\begin{proposition}[Monotonicity]\label{prop:mono14}
Under the intersection update, if $K'\subseteq K$ and both give a nonempty intersection with $C_f$, then the reading set of $K'\cap C_f$ is contained in that of $K\cap C_f$. Knowledge that grows can only narrow what an edition leaves open.
\end{proposition}

The proof is that a subset of a set has the readings of a subset. A check over all $680$ consistent pairs of a faithful edition with two of the later films found no counterexample. The consequence is worth stating for the question the chapter was asked. Within this model, a viewer who has seen more before an edition can find no reading in it that a viewer who has seen less could not. Everything that a later release gains over an earlier one has to come from somewhere the model does not yet contain.

\section{Order and the Path}\label{sec:path14}

Take one of the six editions that leaves all four readings open, together with two later films, the noticing and the noticing and concealing. Three films can be seen in six orders. \Cref{tab:order14} gives, for each order, the openness at each film's first viewing, taking each film's reading set against what the viewer knew when it was seen.

\begin{table}[ht]
\centering
\begin{tabular}{@{}lcc@{}}
\toprule
Order & Openness at each first viewing & Total\\
\midrule
edition, noticed, concealed & 4, 2, 1 & 7\\
edition, concealed, noticed & 4, 1, 1 & 6\\
noticed, edition, concealed & 2, 2, 1 & 5\\
noticed, concealed, edition & 2, 1, 1 & 4\\
concealed, edition, noticed & 1, 1, 1 & 3\\
concealed, noticed, edition & 1, 1, 1 & 3\\
\bottomrule
\end{tabular}
\caption{Openness at the first viewing of each film in each of the six orders, for one fully open edition and the two later films, under the intersection update.}\label{tab:order14}
\end{table}

The final knowledge is the same in every row, and the openness at first viewing is not. The last two orders give the same openness sequence, so openness does not separate every order, and whether the states visited differ in those two is a separate question that the table does not answer. The edition seen first is seen with four readings open, and seen after the concealment film it is seen with one. The total over the three viewings ranges from three to seven. Over all $32$ faithful editions the spread of this total across the six orders is one for ten of them, two for fourteen, three for two, and four for six, which are the fully open ones, and the order that most often gives the largest total, in $24$ of the $32$ with ties counted, is the edition first, then the noticing, then the concealing. Order does not change what the viewer eventually knows, and it changes how long the question stays open, which is the part of a story that a viewer cannot repeat.

\section{Contradiction and the Latest Film}\label{sec:contra14}

When a later film contradicts the earlier, order does change the final state, and \Cref{prop:order14} says exactly when. Of the $32$ faithful editions, eight commit to the third bolt loose, and a film showing it tight contradicts them. For those eight and that film the revision update gives the film's set when the film is last and the edition's set when it is first, and the final state depends on the order for those eight and for none of the other twenty-four. For the films that are consistent with every faithful edition, the noticing, the concealing, and the third bolt loose, the final state depends on the order for none of the thirty-two.

This is a small result and it fixes what the unreliable narrator costs in the model. A later film that contradicts an earlier one does not narrow the earlier reading. It replaces it, and the replacement stands only while no film follows it. The revision update gives the latest film priority and nothing in the model says that viewers do so.

\section{Repertoires as Structures}\label{sec:struct14}

If order matters to the path, a repertoire is a sequence and not a set. The space of repertoires of $n$ films is then larger by a factor of $n!$: $120$ orders of five films, $3{,}628{,}800$ of ten, and a number with $119$ digits for the eighty productions of Chapter~13. Release dependencies cut the space. A repertoire can be a partially ordered set, where one film is to be seen before another because the second presupposes the first, and the viewings that respect it are its linear extensions.

\begin{table}[ht]
\centering
\begin{tabular}{@{}lr@{}}
\toprule
Five films, with the constraints & Admissible orders\\
\midrule
none & 120\\
the fourth film after the third (a reveal after its setup) & 60\\
a chain of three, the others free & 20\\
a chain of three and a chain of two & 10\\
a chain of five & 1\\
\bottomrule
\end{tabular}
\caption{Number of orders of five films that respect a set of precedence constraints.}\label{tab:ext14}
\end{table}

Each precedence halves or divides the orders, and a full chain removes the choice. A series built as one long chain offers a viewer one order, and a series built as independent episodes offers all of them. This is the structural trade of a repertoire: the more the films depend on each other, the more a later film can change an earlier one, and the fewer orders a viewer can choose among.

\section{The Tenth Year and the First}\label{sec:tenth14}

The question the outline poses is whether releasing a film to a viewer in the tenth year of life does something that releasing it in the first would miss. In the model of this chapter the viewer's knowledge at a given age is the intersection of the films seen before, so it only narrows with age, and \Cref{prop:mono14} says the readings of an edition only narrow with it. A release in the first year sees at least as many readings open as one in the tenth. The model therefore answers that a later release can only lose openness, and what it gains must come from outside it.

Three possible gains are left as questions. A reading may require the viewer to know a premise, so that a viewer who lacks it cannot recognize the reading at all even though the histories that support it are not excluded. If so, the set of readings a viewer can recognize is not monotone in knowledge, and a later release can recognize what an earlier one could not. The premise of a reading, and whether viewers with and without it recognize the reading differently, is a measurement and the model contains no value for it. A second gain is that the viewer of the tenth year has different tolerances, those of Chapter~2, so that the same edition and a different viewer give a different equivalence. A third is that the first viewing also sets up the later, since a viewer who saw the film early has a memory of it that the later viewer lacks, and the memory is itself a film that the later release is read against. The model has no memory, because its viewer holds only a set of histories.

An experiment that would settle the first question would hold the edition fixed, vary the viewer's prior films, and ask which readings viewers report, for readings that have a premise and for readings that do not. If the second set is monotone in the prior knowledge and the first is not, the tenth year does something that the first misses. If both are monotone, it does not, and the release date is a matter of what the viewer wants to be told and when. The chapter does not make the prediction.

\section{What the Chapter Establishes and What It Assumes}\label{sec:limits14}

The propositions are theorems of the model: that order matters to the final state only through contradiction, that reading sets narrow with knowledge, and that the openness at first viewing can differ with the order. The counts are on the engine scene, a story of twenty-two histories with four readings, and on five later films chosen by what they commit to. The two updates are the simplest available and neither is a model of understanding. A viewer here holds a set of histories and has no memory, no attention, and no forgetting, and a film commits to a set of histories exactly when its shots agree with them, which Chapter~3 defined and did not measure. The statement that a viewer's experience is a path of openness values is itself a modeling choice, since a viewer may not experience openness as the size of a set of readings.

The remaining chapters leave the single viewer and the toy scene. Chapter~15 asks what a production must exclude from its editions, and returns to counting.

The script \texttt{verify\_ch14.py} reproduces each figure of this chapter.

\begin{thebibliography}{99}

\bibitem{flyxion2026gf}
Flyxion.
\newblock \emph{Generative Film: Categories, Sheaves, and Types for Cinema}.
\newblock Unpublished monograph, October 2026.

\end{thebibliography}

\end{document}
